% MCLF Book
% Models of Curves and Valuations
%
% Copyright (c) 2026 Stefan Wewers and contributors.
%
% License: CC BY 4.0. See the file LICENSE in the repository root.
%
% This file is part of the open book project hosted at:
% https://github.com/MCLF/mclf-book
%
% The text is work in progress. For citation, please refer to a
% specific Git commit or release.

\documentclass[11pt,a4paper,oneside]{book}

% MCLF Book
% Models of Curves and Valuations
%
% Copyright (c) 2026 Stefan Wewers and contributors.
%
% License: CC BY 4.0. See the file LICENSE in the repository root.
%
% This file is part of the open book project hosted at:
% https://github.com/MCLF/mclf-book
%
% The text is work in progress. For citation, please refer to a
% specific Git commit or release.



% Basic input and font setup.
% With a modern LaTeX installation, utf8 is usually the default,
% but keeping this explicit is harmless for pdfLaTeX.
\usepackage[utf8]{inputenc}
\usepackage[T1]{fontenc}

% Standard mathematical packages.
\usepackage{amsmath}
\usepackage{amssymb}
\usepackage{amsthm}
\usepackage{mathtools}
\usepackage{stix}
%\usepackage{unicode-math}

% Configurable labels for standard list environments.
\usepackage{enumitem}

% Commutative diagrams.
\usepackage{tikz-cd}
\newcommand{\cdpunkt}{\mathrlap{\, .}}
\tikzset{
    symbol/.style={
        draw=none,
        every to/.append style={
            edge node={node [sloped, allow upside down, auto=false]{$#1$}}}
    }
}
\usepackage{tikzit}
\input{basic.tikzstyles}

\usepackage{comment}

% Graphics.
\usepackage{graphicx}
\usepackage{float}
\usepackage{tcolorbox}

% Visible notes for mathematical issues that require human resolution.
\newtcolorbox{authorwarning}{
    colback=red!5!white,
    colframe=red!75!black,
    fonttitle=\bfseries,
    title=Author input required
}

% References and hyperlinks.
% Load hyperref near the end.
\usepackage[
    colorlinks=true,
    linkcolor=blue,
    citecolor=blue,
    urlcolor=blue
]{hyperref}

% Theorem environments.
\theoremstyle{plain}
\newtheorem{theorem}{Theorem}[chapter]
\newtheorem{proposition}[theorem]{Proposition}
\newtheorem{lemma}[theorem]{Lemma}
\newtheorem{corollary}[theorem]{Corollary}

\theoremstyle{definition}
\newtheorem{definition}[theorem]{Definition}
\newtheorem{example}[theorem]{Example}
\newtheorem{construction}[theorem]{Construction}
\newtheorem{assumption}[theorem]{Assumption}

\theoremstyle{remark}
\newtheorem{remark}[theorem]{Remark}
\newtheorem{remarks}[theorem]{Remarks}

% Number equations independently by chapter.
\numberwithin{equation}{chapter}

% Slightly more flexible spacing in tables.
\renewcommand{\arraystretch}{1.1}

\input{macros}

\title{MCLF - the book\\[0.5em]
\large Models of curves over local fields}
\author{Edited by Stefan Wewers}
\date{\today}

\begin{document}

\frontmatter

\maketitle

\tableofcontents

\chapter*{Preface}
\addcontentsline{toc}{chapter}{Preface}
\markboth{Preface}{Preface}

The MCLF project (\emph{Models of Curves over Local Fields}) develops a SageMath toolbox for computations with models of curves over local fields, with particular emphasis on semistable reduction.
The project is hosted at \url{https://github.com/MCLF/mclf}.

\medskip

This book grew out of an Oberseminar held during the winter semester 2024/25 at the Institute of Algebra and Number Theory at Ulm University.
The seminar was organized as a short series of talks under the title \emph{Models of curves}.
Its original aim was to develop, from both a theoretical and a computational perspective, some of the basic ideas underlying the MCLF project.

\medskip\noindent
The seminar consisted of five talks:

\begin{itemize}
    \item Miriam Ní Chobhtaigh: \emph{Models of curves and valuations};
    \item Max Schwegele: \emph{Proof of Theorem 1};
    \item Xiaodong Zhang: \emph{Models with reduced special fiber}; % notation-check: allow quoted title
    \item Robert Nowak: \emph{The Epp problem for \(p\)-cyclic covers of curves};
    \item Kletus Stern: \emph{Mac Lane valuations}.
\end{itemize}


The initial plan was simply to revise and combine the notes prepared for these talks into a larger set of seminar notes.
As work on the project progressed, however, it became clear that the material would benefit from a broader and more systematic treatment.
In particular, the relation between algebraic models, valuations, Berkovich geometry, Mac Lane theory, and explicit computations suggested a structure extending well beyond the original five talks.

The present book is the result of this development. The original seminar notes were compiled, edited and extended by Maximilian Fülle. The content of the first four talks is the germ of Part I, \emph{Models and valuations}. The notes for the fifth talk form the starting point of Part II, \emph{Analytic and computational tools}.

The book is a collaborative and evolving project.
Individual chapters identify their authors, while the volume as a whole is edited by Stefan Wewers.
The chapter attributions refer to the finished chapters, not merely to the speakers in the original seminar: in some cases the seminar notes have been substantially revised, reorganized, or supplemented by additional authors.

We thank all participants in the Oberseminar for their talks, notes, questions, and discussions.
Without the seminar and the work invested in its preparation, this book would not have been started.

\chapter*{Notation and conventions}
\addcontentsline{toc}{chapter}{Notation and conventions}
\markboth{Notation and conventions}{Notation and conventions}

\paragraph*{Valued fields.}
We write $(K,v_K)$ for a valued field, with valuations written additively, and use
\[
    \OK, \qquad \mK, \qquad k=\OK/\mK, \qquad
    \Gamma_K=v_K(K^\times)
\]
for its valuation ring, maximal ideal, residue field, and value group.
When a discrete valuation is normalized, our convention is $v_K(K^\times)=\ZZ$.
Completeness, discreteness, and assumptions on the residue field are stated where they are needed and are not imposed globally unless explicitly indicated.

\paragraph*{Curves and function fields.}
The letter $X$ usually denotes a smooth, projective, absolutely irreducible curve over $K$, called a \emph{nice curve} in Chapter~\ref{ch:models-of-curves-over-valued-fields}.
Its function field is denoted by $F_X=K(X)$.

\paragraph*{Models and fibres.}
A model of $X$ over $\OK$ is generally denoted by $\X$.
If $\eta$ and $s$ are respectively the generic and closed points of $\Spec\OK$, then
\[
    \X_\eta \quad\text{and}\quad \X_s
\]
denote the generic and special fibres.
Thus the book uses the French-style notation $\X_\eta$, not $\X_g$, for the generic fibre. % notation-check: allow documented deprecated form
For an irreducible component $Z$ of a special fibre, its generic point is denoted by $\xi_Z$, or simply $\xi$ when no confusion can arise; the associated geometric valuation is $v_Z$, and the finite set of component valuations is $V(\X)$.

\paragraph*{Base change.}
For a scheme $Y$ over a ring $R$ and a ring homomorphism $R\to S$, we write
\[
    Y_S \coloneqq Y\times_{\Spec R}\Spec S.
\]
After this definition, we use the subscript notation $Y_S$; in particular, base changes of $X$ and $\X$ are denoted by $X_L$ and $\X_{\OL}$.
For rings, modules, algebras, and sheaves, base change is denoted by $\otimes_R$.
When $L/K$ is finite, the normalized base change of $\X$ is denoted by
\[
    (\X_{\OL})^\sim.
\]
Function-field composita are written $F_X\mathbin{\cdot}L$.


\mainmatter

\part[Models and valuations]{Models and valuations\\[1.5em]
\large\normalfont
Miriam Ní Chobhtaigh, Max Schwegele,\\
Xiaodong Zhang, Kletus Stern, Robert Nowak,\\
Stefan Wewers, and Maximilian Fülle\\[2em]
\begin{minipage}{0.78\textwidth}
\normalfont\normalsize\raggedright
This part develops the algebraic foundations and valuation theory used throughout the book. It introduces models of curves, attaches geometric valuations to components of their special fibres, establishes the correspondence between models and finite sets of these valuations, and studies their behaviour under finite extensions of the base field.
\end{minipage}}
\label{part:models-and-valuations}

%! Author = Maximilian Fülle
%! Date = 29.05.2026


\chapter{Models of curves over valued fields}\label{ch:models-of-curves-over-valued-fields}

This chapter introduces models of smooth projective curves over valued fields, together with their morphisms and the resulting partial order by domination.
After discussing normality and basic examples, we attach a geometric valuation and a multiplicity to every irreducible component of the special fibre; the resulting finite set of valuations is the starting point for the correspondence developed in Chapter~\ref{ch:correspondence-between-models-and-valuations}.

\section{Valued fields, curves, and models}\label{sec:valued-fields-curves-and-models}
We begin by fixing the valued-field notation and standing hypotheses used throughout Part~\ref{part:models-and-valuations}, and then introduce the basic notion of a model of a curve.

Let $K$ be a field endowed with an additive valuation of rank $1$, $v_K: K \to \hat{\RR} \coloneqq \RR\cup \lbrace\infty\rbrace$.
Let $\Gamma_K \coloneqq v_K(K^*) \subset \RR$ be the value group of $v_K$, $\OK \subset K$ its valuation ring with
maximal ideal $\mK$ and $k \coloneqq \OK /\mK$ its residue field.

For most of the first chapter, we make the following additional assumptions:
\begin{assumption} \label{ass:v_K}
    \begin{enumerate}[label=\textup{(\alph*)}]
        \item $v_K$ is discrete, i.e.\ $\Gamma_K$ is a discrete subgroup of $\RR$,
        \item $K$ is complete with respect to $v_K$,
        \item $k$ is perfect.
    \end{enumerate}
\end{assumption}

Throughout the monograph, $X$ will denote a \textit{nice} curve over $K$,
that is a smooth, projective and absolutely irreducible curve over $K$.

\begin{definition} \label{def:model}
    A \textit{model} $\X$ of $X$ is a proper, flat and normal $\OK $-scheme
    with an identification of its generic fibre $\X_\eta$ with $X$.
\end{definition}

\begin{example}
    Let $K = \QQ_p$, $\OK  = \ZZ_p$ and $X = \PP^1_K$.
    Then $\X = \PP^1_{\OK}$ is a model of $X$ with the evident identification
    $\X_\eta = \X_K \cong \PP^1_K$.
    The special fibre $\X_s$ has one irreducible component which is isomorphic to $\PP^1_k$.
\end{example}

\begin{remark} \label{rem:flatIffIntegral}
    The condition of a model being flat is equivalent to it being integral:
    if $\X$ is integral and its structural map is compatible with the given generic fibre, then every local ring $\OO_{\X,x}$ is a domain containing $\OK$.
    Thus $\OO_{\X,x}$ is a torsion-free module over the discrete valuation ring $\OK$, hence flat by~\cite[Corollary 1.2.14]{Liu}.
    Conversely, the integrality of a model $\X$ of $X$ follows from its flatness and the integrality of $X$ by~\cite[Proposition 4.3.8]{Liu}.
\end{remark}


\section{The category of models and modifications}\label{sec:category-of-models-and-modifications}
We now introduce the morphisms used to compare models of a fixed curve.
They turn the collection of models into a partially ordered system in which any two models admit a common dominating model.

\begin{definition}\label{def:modification-of-models}
    Let $\X$ and $\X'$ be models of $X$.
    A \textit{modification} $\varphi\colon\X'\to\X$ is a morphism of $\OK$-schemes whose restriction to the generic fibres is the identity on $X$ under the fixed identifications $\X'_\eta\cong X\cong\X_\eta$.
    In this situation, we say that $\X'$ \textit{dominates} $\X$.
\end{definition}

\begin{lemma}\label{lem:uniqueness-of-modifications}
    Between two fixed models of $X$, there is at most one modification.
\end{lemma}

\begin{proof}
    Let $\varphi,\psi\colon\X'\to\X$ be modifications.
    They agree on the generic fibre $\X'_\eta$, which is dense in the integral scheme $\X'$.
    Since $\X$ is separated, the equalizer of $\varphi$ and $\psi$ is closed in $\X'$.
    It contains the dense generic fibre and is therefore all of $\X'$, so $\varphi=\psi$.
\end{proof}

\begin{proposition}\label{prop:elementary-properties-of-modifications}
    Let $\varphi\colon\X'\to\X$ be a modification of models.
    Then $\varphi$ is proper, dominant, and surjective.
    Moreover, there is a finite set $S\subset\X_s$ of closed points such that $\varphi$ is an isomorphism over $\X\setminus S$ and, for every $x\in S$, the fibre $\varphi^{-1}(x)$ is a connected projective curve over $\kappa(x)$.
\end{proposition}

\begin{proof}
    Let $\varphi\colon\X'\to\X$ be a modification.
    Since $\X$ is separated over $\Spec\OK$, the graph morphism
    \[
        \Gamma_\varphi\colon\X'\longrightarrow\X'\times_{\Spec\OK}\X
    \]
    is a closed immersion.
    The projection $\X'\times_{\Spec\OK}\X\to\X$ is proper because it is the base change of the proper morphism $\X'\to\Spec\OK$.
    Their composite is $\varphi$, so $\varphi$ is proper.
    Its image contains the generic fibre $X$, which is dense in $\X$, and hence $\varphi$ is dominant.
    Finally, a proper morphism has closed image, so the image of $\varphi$ is both dense and closed in $\X$ and must equal $\X$.

    A proper birational morphism to a normal scheme is an isomorphism at every codimension-one point of the target; this follows by applying the valuative criterion of properness to the corresponding discrete valuation ring.
    Hence the complement of the maximal open subset of $\X$ over which $\varphi$ is an isomorphism is a closed subset of codimension at least two.
    Since $\X$ is a two-dimensional Noetherian scheme and $\varphi$ is already an isomorphism on the generic fibre, this complement is a finite set $S\subset\X_s$ of closed points.

    By the connectedness form of Zariski's Main Theorem, the fibres of $\varphi$ are connected; see~\cite[\href{https://stacks.math.columbia.edu/tag/0AY8}{Lemma 0AY8}]{Stacks}.
    If $x\in S$, the fibre over $x$ cannot be zero-dimensional: otherwise $\varphi$ would be finite in a neighbourhood of $x$, and a finite birational morphism to the normal scheme $\X$ would be an isomorphism there.
    Thus $\varphi^{-1}(x)$ is a proper one-dimensional scheme over $\kappa(x)$, hence a projective curve.
\end{proof}

We write
\[
    \X\leq\X'
\]
if there is a modification $\X'\to\X$; thus the larger model is the one that dominates the smaller model.
Here and below, models that are isomorphic as models of $X$ are identified when discussing this relation.

\begin{proposition}\label{prop:partial-order-of-models}
    The relation $\leq$ is a partial order on the set of isomorphism classes of models of $X$.
    The category of models and modifications is equivalent to the partially ordered set regarded as a category.
\end{proposition}

\begin{proof}
    The identity morphism gives reflexivity, and composition of modifications gives transitivity.
    If $\X\leq\X'$ and $\X'\leq\X$, there are modifications $\varphi\colon\X'\to\X$ and $\psi\colon\X\to\X'$.
    Both composites are modifications and hence, by Lemma~\ref{lem:uniqueness-of-modifications}, are the respective identity morphisms.
    Thus $\varphi$ and $\psi$ are mutually inverse, which proves antisymmetry on isomorphism classes.
    The same uniqueness lemma says that there is at most one morphism between any two models, so choosing one representative of each isomorphism class gives the asserted equivalence of categories.
\end{proof}

\begin{proposition}\label{prop:common-dominating-model}
    Any two models $\X_1$ and $\X_2$ of $X$ admit a common dominating model $\X_3$.
\end{proposition}

\begin{proof}
    Consider the diagonal morphism determined by the fixed identifications of the generic fibres,
    \[
        X\longrightarrow\X_1\times_{\Spec\OK}\X_2.
    \]
    Let $\mathcal Z$ be the schematic closure of its image.
    The product $\X_1\times_{\Spec\OK}\X_2$ is proper over $\Spec\OK$, and hence so is its closed subscheme $\mathcal Z$.
    Moreover, $\mathcal Z$ is integral and its generic fibre is the diagonal copy of $X$.
    It is therefore flat over the discrete valuation ring $\OK$ by Remark~\ref{rem:flatIffIntegral}.

    Let $\X_3$ be the normalization of $\mathcal Z$ in $F_X$.
    The completeness assumption in Assumption~\ref{ass:v_K} enters here: it implies that $\OK$ is excellent, so the normalization is finite; see~\cite[Chapter 8, in particular Theorem 8.2.39 (d)]{Liu}.
    Consequently, $\X_3$ is proper over $\Spec\OK$.
    It is integral and hence flat over $\OK$, and it is normal by construction.
    Since $X$ is normal, normalization does not change the generic fibre, so $(\X_3)_\eta=X$.
    Thus $\X_3$ is a model of $X$.
    Composing the normalization map with the two projections gives modifications $\X_3\to\X_i$ for $i=1,2$.
\end{proof}

\begin{remark}\label{rem:models-form-cofiltered-category}
    Proposition~\ref{prop:common-dominating-model} says that the partially ordered set of models is directed upward.
    Since modifications point from a dominating model to a dominated model, the category of models is correspondingly cofiltered and may be viewed as an inverse system.
    Common dominating models may also be produced through suitable blowups followed by normalization.
\end{remark}

Chapter~\ref{ch:correspondence-between-models-and-valuations} will identify this order with inclusion between the corresponding finite sets of geometric valuations.


\section{Normality and examples of models}\label{sec:normality-and-examples-of-models}
To recognize when a proper flat scheme is a model, it remains to understand its normality. We first recall a criterion in terms of the special fibre and then apply it to several examples.

\begin{definition}
    A \textit{pre-model} of $X$ is a flat and proper $\OK$-scheme $\X$ with $\X_\eta = X$.
    That is, $\X$ satisfies all the conditions of a model except that it need not be normal.
\end{definition}

The following normality criterion is a reformulation of~\cite[Lemma 2.1]{ArzdorfWewers} in the present notation.

\begin{lemma}\label{lem:conditions-for-normality}
    Let $\X$ be a pre-model of $X$.
    Then $\X$ is normal (i.e.\ a model) if and only if the following two conditions hold.
    \begin{enumerate}[label=\textup{(\alph*)}]
        \item The special fibre $\X_s$ has no embedded points, i.e.\ no point $x\in \X_s$ \label{lem:conditions-for-normality-i}
            satisfying $\m_x = \Ann{R}{r}$ for some $r\in R \coloneqq \OO_{\X_s,x}$ (see~\cite[Definition 7.1.6]{Liu}).
        \item The local ring $\OO_{\X,\xi}$ is a discrete valuation ring for every generic point $\xi \in \X_s$. \label{lem:conditions-for-normality-ii}
    \end{enumerate}
\end{lemma}

\begin{proof}
    By Serre's criterion for normality~\cite[Theorem 8.2.23]{Liu},
    we need to check the following properties:
    \begin{itemize}
        \item[$(\mathrm{R}_1)$] $\X$ is regular at $x$ when $\codim{x} \leq 1$, and
        \item[$(\mathrm{S}_2)$] $\depth{\OO_{\X,x}} \geq 2$
    \end{itemize}
    for all $x\in \X$. \\
    The points on $\X$ with codimension $1$ are precisely the points on the generic fibre, which is regular by assumption,
    and the generic point on the special fibre.
    The codimension-one local rings under consideration are one-dimensional Noetherian local domains.
    Such a ring is regular if and only if it is a discrete valuation ring, so $(\mathrm{R}_1)$ is equivalent to
    \ref{lem:conditions-for-normality-ii}.
    By~\cite[Proposition 8.2.11]{Liu}, condition $(\mathrm{S}_2)$ is equivalent to the condition
    that $\depth{\OO_{\X_s, x}} \geq 1$ for all $x \in \X_s$.
    This condition always holds at a generic point, and at any other point $x$ it is equivalent to the condition that $x$ is not an embedded point.
\end{proof}

\begin{example}
    \begin{enumerate}[label=\textup{(\alph*)}]
        \item Consider a pre-model with special fibre
            \begin{align}
                \X_s \coloneqq \Spec{k[x,y]/(x^2,xy)},
            \end{align}
            which is topologically equivalent to $\AA^1_k$.
            The point corresponding to the maximal ideal $(x,y)$ is an embedded point.
        \item We consider the pre-model of a smooth curve with an embedded point on the special fibre described in \cite[Example \textrm{III}.9.8.4]{Hartshorne}.
            Let $C$ be the curve in $\PP^3_{\QQ_p} = \Proj{\QQ_p[X,Y,Z,W]}$, which is given by the ideal
            \begin{align}
                I = \left(p^2(x+1) - z^2, px(x+1) - yz, xz - py, y^2 - x^2(x+1) \right)
            \end{align}
            on the affine chart $W \neq 0$.
            $p$ is not a zero divisor in $\ZZ_p[x,y,z]/I$.
            Thus, the projective $\ZZ_p$-scheme $\mathcal{C}$ defined by $I$ is flat and hence a pre-model of $C$.
            On the affine chart $W \neq 0$, the special fibre $\mathcal{C}_s$ is given by the ideal
            \begin{align}
                I_p = (z^2, yz, xz, y^2 - x^2(x+1))
            \end{align}
            and has support equal to the nodal cubic curve $y^2 = x^2(x+1)$.
            At any point where $x \neq 0$, $\mathcal{C}_s$ is reduced.
            However, in the stalk at $(0,0,0)$, the element $z$ is a nonzero nilpotent because $z^2=0$.
            Thus, $(0,0,0)$ is an embedded point, and $\mathcal{C}$ is not a model of $C$ by Lemma \ref{lem:conditions-for-normality}.
    \end{enumerate}
\end{example}


\begin{comment}
    \begin{example}
    In this example, we consider models of curves in the plane $\PP^2_K$.
    Such a (smooth, projective and irreducible) curve $X$ is defined by a homogeneous irreducible polynomial $F(X,Y,Z) \in K[X,Y,Z]$.
    We show that if $F(X,Y,Z)$ is integral (that is all coefficients of $F$ are in $\OK$) and primitive, a model of $X$ is given by
    \begin{align}
        \X = \Proj{\OK[X,Y,Z]/(F)}.
    \end{align}
    Is is clearly a proper $\OK$-scheme.
    For the flatness, it suffices to  show that $R \coloneqq \OK[X,Y,Z]/(F(X,Y,Z))$ is a torsion-free $\OK$-module by Remark \ref{rem:flatIffIntegral}.
    This is the case if and only if $\bar{F} \neq 0 \in k[X,Y,Z]$, which is equivalent to $F$ being primitive.
    For the normality, we use the following argument which is similar to Example 8.2.26 of \cite{Liu}.
    Since $\X$ is a Cohen-Macaulay scheme, it is normal if and only if it is normal at the points of codimension $1$.
    At this points, normality is equivalent to regularity.
    These are exactly the closed points on the generic fibre $\X_\eta=X$ and the generic points of the special fibre $\X_s$.
    Let $\xi \in \X$ be a generic point of the special fibre $\X_s$ and let $\pi\in \OK$ be a uniformising parameter.
    On the affine chart $U: Z \neq 0$ with $x = X/Z$ and $y=Y/Z$, the point $\xi$ corresponds to the ideal $\m$ that is generated by $\pi$ and a
    polynomial $G(x,y) \in \OK[x,y]$ whose image $\bar{G}(x,y) \in k[x,y]$ is irreducible.
    We can write
    \begin{align}
        F(x,y) = G(x,y)^r H_1(x,y) + \pi^e H_2(x,y),
    \end{align}
    for $H_i(x,y)\in \OK[x,y]$ and $r,e \geq 1$ and with $\bar{H}_1(x,y) \notin \bar{G}(x,y)k[x,y]$ and $\bar{H}_2(x,y) \neq 0$.
    By Corollary 4.2.12 of \cite{Liu}, $\xi$ is regular if $F \notin \m^2$.
    Thus, $\X$ is normal at $\xi$ if and only if either $r=1$; or $e=1$ and $\bar{H}_2(x,y) \neq 0$.
    \end{example}
\end{comment}

\begin{example}\label{ex:model-of-plane-curve}
    We discuss the important example of plane curves.
    A smooth and absolutely irreducible curve in $\PP^2_K$ is defined by a homogeneous irreducible polynomial $F \in K [X,Y,Z]$.
    After multiplying $F$ by a suitable element of $K^\times$, we may assume that $F$ is integral (i.e., all coefficients are in $\OK$) and primitive.
    Then the $\OK$-scheme
    \begin{align}
        \X = \Proj{\OK[X,Y,Z]/(F)}
    \end{align}
    is an obvious candidate for a model of $X$.
    $\X$ is proper as a closed subscheme of $\PP^2_K$ and flat by Remark \ref{rem:flatIffIntegral} since
    $\OK[X,Y,Z]/(F)$ is torsion-free by irreducibility and primitivity of $F$.
    Hence, $\X$ is a model if and only if it is normal.
    We rephrase the discussion of normality in \cite[Example 8.2.26]{Liu} for the present setting.
    Since $\X$ is Cohen--Macaulay, normality of $\X$ is equivalent to normality at the points of codimension $1$,
    which in turn is equivalent to regularity at these points.
    The points of codimension $1$ are the closed points of the (smooth) generic fibre $\X_\eta = X$,
    and the generic points of the special fibre $\X_s$.
    Now let $\xi \in \X_s$ be a generic point.
    On the affine chart $Z \neq 0$, with $x = X/Z$ and $y = Y/Z$, the point $\xi$ corresponds to the ideal $\m$ generated by a uniformising parameter $\pi \in \OK$
    and a polynomial $G \in K[x,y]$ whose reduction $\bar{G} \in k[x,y]$ is irreducible.
    We can write
    \begin{align}
        F(x,y) = G(x,y)^r H_1(x,y) + \pi^e H_2(x,y),
    \end{align}
    where $\bar{H}_1(x,y) \notin \bar{G}(x,y)k[x,y]$ and $\bar{H}_2(x,y)\neq 0$.
    By \cite[Corollary 4.2.12]{Liu}, $\X$ is regular at $\xi$ if and only if $F \notin \m^2$,
    which in turn is the case if and only if either $r=1$, or $e=1$ and $\bar{H}_2 \notin \bar{G}(x,y)k[x,y]$.

\end{example}

\begin{example}\label{ex:particular-model-of-plane-curve}
    \begin{enumerate}[label=\textup{(\alph*)}]
        \item Consider the curve $X$ defined by $y^2 = x^3 + px$ over $\QQ_p$ with $p \neq 2$.
            We have seen that $\X = \Proj{\OK[X,Y,Z]/(X^3 + pXZ^2 - Y^2Z)}$ is a proper and flat $\OK$-scheme with $\X_\eta = X$.
            The special fibre is given by $\X_s = \Proj{k[X,Y,Z]/(X^3 - Y^2Z)}$.
            On the affine cover $Z\neq 0$, its generic point $\xi$ corresponds to the ideal $\m = (p, x^3-y^2)$.
            Thus, in the notation of Example \ref{ex:model-of-plane-curve}, $r=e=1$; that is, $\X$ is normal and therefore a model of $X$.
        \item Let $X: y^3 = 1 + x^3 + 3x$ be a curve over $\QQ_3$.
            The scheme
            \begin{align}
                \X = \Proj{\OK[X,Y,Z]/(Y^3 - Z^3 - X^3 - 3XZ^2)}
            \end{align}
            has a non-reduced special fibre:
            On the affine chart $Z \neq 0$ we have $\X_s: 0 = y^3 - x^3 - 1 = (y - x - 1)^3$ over $\FF_3$.
            However, $\X$ is normal and thus a model of $X$ since
            \begin{align*}
                F(x,y) &= G(x,y)^r H_1(x,y) + 3 H_2(x,y) \\
                    &\coloneqq (y-x-1)^3 + 3(-x^2y + xy^2 + x^2 - 2xy + y^2 - y),
            \end{align*}
            i.e. $r=3$ but $e=1$ and $\bar{H}_2(x,y) \notin \bar{G}(x,y)\FF_3[x,y]$.
    \end{enumerate}
\end{example}

% Model with reduced special fibre: X: 0 = (y-1)(x^3+2) +3x

\section{Components and geometric valuations}\label{sec:components-and-geometric-valuations}
We now associate a valuation of the function field to each irreducible component of the special fibre. This construction also relates the multiplicity of a component to the value group of its valuation.

In this section, we assume the valuation $v_K$ to be \textit{normalized}, i.e. $v_K(K^\times) = \ZZ$.

\begin{definition}
    Let $F_X$ be the function field of $X$ over $K$.
    A valuation $v: F_X \to \hat{\RR}$ is said to be \textit{of type \textrm{II}}, if
    \begin{enumerate}[label=\textup{(\alph*)}]
        \item $v|_K = v_K$,
        \item $v$ is discrete, and
        \item the extension of residue fields $k(v)|k$ has transcendence degree $1$.
    \end{enumerate}
    A valuation of this type is also known as a \textit{geometric valuation}.
    We write $V_\textrm{II}(F_X)$ for the set of all valuations of type \textrm{II} of $F_X$.
    The number
    \begin{align}
        m(v) \coloneqq [\Gamma_v : \Gamma_K]
    \end{align}
    is called the \textit{multiplicity of $v$}, where $\Gamma_v \coloneqq v(F_X^*)$ is the value group.
    We say that $v$ is \textit{weakly unramified} if $m(v) = 1$.
\end{definition}

\begin{example}
    The \textit{Gauß valuation} on the polynomial ring $K[x]$ is defined as
    \begin{align}
        v_{\scriptscriptstyle G}: K[x] &\to \hat{\QQ}, \; v_{\scriptscriptstyle G}\left( \sum_i a_i x^{i} \right) = \min_i(v_K(a_i)).
    \end{align}
    Since $k(v_{\scriptscriptstyle G}) = k(\bar{x})$ (where $\bar{x}$ is the image of $x$ in $k(v_{\scriptscriptstyle G})$), the Gauß valuation
    extends naturally and uniquely to $K(x)$, so $v_{\scriptscriptstyle G} \in V_\textrm{II}(K(x))$.
\end{example}

Let $\X$ be a model of $X$.
We wish to assign a valuation in $V_\textrm{II}(F_X)$ to each generic point of the special fibre $\X_s$.
Suppose that $Z$ is an irreducible component of $\X_s$ with generic point $\xi$.
Since $\xi$ has codimension $1$ in $\X$, the stalk $\OO_{\X,\xi}$ is a discrete valuation ring\@.
Thus, we can define a valuation on $F_X$,
\begin{align} \label{eq:ValuationOfComponent}
    v_Z: \mathrm{Frac}(\OO_{\X, \xi}) = F_X \to \hat{\QQ}.
\end{align}
Since $\X$ is a flat $\OK $-scheme, $\OO_{\X, \xi}$ dominates $\OK$, and we may normalize $v_Z$ so that its restriction to $K$ is $v_K$.
We set
\begin{align}
    V(\X) \coloneqq \lbrace v_Z: \; Z \subset \X_s \text{ is an irreducible component}\rbrace.
\end{align}

\begin{proposition}\label{prop:WellDefOfInducedValuation}
    The valuation defined in~\eqref{eq:ValuationOfComponent} is of type \textrm{II}.
\end{proposition}

\begin{proof}
    As mentioned above, we can assume $v_Z|_K = v_K$.
    The residue field of the valuation is
    \begin{align}
        \kappa(v_Z)=\kappa(\xi)=k(Z).
    \end{align}
    Since $Z$ is an integral curve over $k$, it follows that
    $\trdeg_k \kappa(v_Z)=1$, showing $v_Z \in V_\textrm{II}(F_X)$.
\end{proof}

\begin{definition}
    The \textit{multiplicity} of the irreducible component $Z$ of $\X_s$ containing the generic point $\xi$ is defined as
    \begin{align}
        m(Z)
        \coloneqq \length_{\OO_{\X,\xi}}\bigl(\OO_{\X_s,\xi}\bigr)
        = \length_{\OO_{\X,\xi}}\bigl(\OO_{\X,\xi}/\pi\OO_{\X,\xi}\bigr),
    \end{align}
    where $\pi$ is a uniformising parameter of $\OK$.
\end{definition}

\begin{proposition}\label{prop:multiplicities-coincide}
    We have $m(v_Z) = m(Z)$.
\end{proposition}

\begin{proof}
    Since the model $\X$ is normal, the stalk $\OO_{\X,\xi}$ is a discrete valuation ring.
    Let $\pi_\xi$ be a uniformising parameter of this ring.
    There is a unit
    $u\in\OO_{\X,\xi}^{\times}$ such that
    \begin{align}
        \pi=u\pi_\xi^{m(Z)}.
    \end{align}
    Indeed, the exponent on the right is the length of
    $\OO_{\X,\xi}/\pi\OO_{\X,\xi}=\OO_{\X_s,\xi}$.
    Since $v_Z(\pi)=v_K(\pi)=1$, the displayed relation gives
    $v_Z(\pi_\xi)=1/m(Z)$.
    Hence
    \begin{align}
        m(Z)=[\Gamma_{v_Z}:\Gamma_K]=m(v_Z).
    \end{align}
\end{proof}

The finite set $V(\X)$ motivates the classification of models developed in the next chapter.

\chapter{The correspondence between models and valuations}\label{ch:correspondence-between-models-and-valuations}
\chaptermark{Models--valuations correspondence}

For the proof of Theorem~\ref{thm:bijectionModelsAndComponents}, it remains to show the surjectivity of the map (\eqref{eq:mapFromModelToValuation}),
i.e., for a given finite, non-empty set of valuations
$V \subset V_\textrm{II}(F_X)$, there is a model $\X$ of $X$ such that $V = V(\X)$.
We will prove this in two steps, which we state in the following proposition.

\begin{proposition}[{{\cite[3.14, 3.16]{Rueth}}}, {{\cite[8.3.36]{Liu}}}] \label{prop:constructing_models}
    \begin{enumerate}[(i)]
        %\item[]
        \item Given model $\X$ and a valuation $v \in V_\textrm{II}(F_X)$, there exists a modification
        $\varphi \colon \X' \to \X$ such that $v \in V(\X')$. \label{prop:constructing_models_i}
        \item Given a model $\X$ and a nonempty subset $V \subset V(\X)$,
        there exists a modification $\varphi \colon \X \to \X'$ such that $V = V(\X')$. \label{prop:constructing_models_ii}
    \end{enumerate}
\end{proposition}

After proving Proposition~\ref{prop:constructing_models}, the proof is complete:
Starting with a model $\X$ of $X$, we can apply part (\ref{prop:constructing_models_i}) for every $v \in V \subset V_\textrm{II}(F_X)$
to obtain a new model that induces all valuations in $V$, but possibly more.
We then apply part (\ref{prop:constructing_models_ii}) to construct a model that induces exactly the valuations in $V$.


\section{Proof of the First Part of Proposition~\ref{prop:constructing_models}}\label{sec:proof-of-the-first-part}

We fix a model $\X$ of $X$ and a type II valuation $v \in V_\textrm{II}(F_X)$.

\begin{definition}
    A point $x \in \X$ is called a \textit{centre} of $v$ if there exists a morphism $f \colon \Spec \OO_v  \to \X$ of $\OK $-schemes
    which maps the closed point to $x$ and makes the following diagram commute:
    \[
        \begin{tikzcd}
            \Spec \OO_v \arrow[rr, "f"] &                                 & \X \\
            & \Spec F_X \arrow[ru] \arrow[lu] &
        \end{tikzcd}
    \]
    Here, the left map is induced by the inclusion $\OO_v \hookrightarrow \Frac \OO_v = F_X$,
    and the right map comes from the identification of $X$ with the special fibre $\X_s$.
    Specifically, it is the canonical map $\Spec F_X \cong \Spec \OO_{\X, \eta} \to \X$, where $\eta$ is the generic point of $X$.
\end{definition}

\begin{lemma}
    There exists a unique centre $x \in \X$ of $v$.
\end{lemma}

\begin{proof}
    By the valuative criterion for properness, such a centre exists, and the morphism $f$ is unique
    (see also~\cite[\href{https://stacks.math.columbia.edu/tag/01J5}{Section 01J5, Lemma 01J6}]{Stacks}):
    \[
        % https://tikzcd.yichuanshen.de/#N4Igdg9gJgpgziAXAbVABwnAlgFyxMJZABgBpiBdUkANwEMAbAVxiRAB12BlNGAYwAEAMQD6ADRABfUuky58hFAEZSSqrUYs2nHvwGcA8gZEBpKTJAZseAkTJrq9Zq0QduvQYeM1zs6wqIVSkdNFzcJSXUYKABzeCJQADMAJwgAWyQAJmocCCQlaSTUjMQAZhy8xAKLFPSkMhBcrMKQWpKGprKWtqyKpHKQBiwwMKgIHBxokBDnNkSpCkkgA
        \begin{tikzcd}[column sep=large, row sep=large]
            \Spec F_X \arrow[d] \arrow[r]                 & \X \arrow[d] \\
            \Spec \OO_v \arrow[r] \arrow[ru, "f", dotted] & \Spec \OK  \cdpunkt
        \end{tikzcd}
    \]
\end{proof}

\begin{lemma}[{{\cite[3.11]{Rueth}}}]\label{lem:codim1}
    We have $v \in V(\X)$ if and only if the centre $x$ of $v$ is of codimension one.
    In that case $x$ is the generic point on the special fibre which induces $v$.
\end{lemma}

\begin{proof}
    Assume $v \in V(\X)$.
    Then $v$ is induced by a generic point $\xi \in \X_s$ on the special fibre.
    The canonical morphism $\Spec \OO_{\X, \xi} \to \X$ shows that $v$ has centre $\xi$,  which is indeed a point of codimension one.
    Conversely, suppose $v \in V_\textrm{II}(F_X)$ has center $x \in \X$ of codimension one.
    The morphism $\Spec \OO_v \to \X$ uniquely factors through $\Spec \OO_v \to \Spec \OO_{\X, x}$, so $\OO_v$ dominates $\OO_{\X, x}$.
    Since these are both valuation rings with common fraction field $K$, it follows that they are in fact equal.
    Thus \(v \in V(\X)\).
\end{proof}

\begin{lemma}[{{\cite[3.12, 3.13]{Rueth}}}] \label{lem:condition_induced}
    Let $x \in \X$ be the centre of $v$, and let $f \in \OO_v^* \subset F_X$ be such that the reduction $\bar{f} \in k(v)$ is transcendental over $k$.
    Then:
    \begin{enumerate}[(a)]
    \item If $f \in \OO_{\X,x}$, then $v \in V(\X)$ \label{lem:condition_induced_i}
    \item If the rational map $\X \dashrightarrow \PP^1(\OK )$ induced by $f$ is defined at the center $x$ of $v$,
    then $v \in V(\X)$. \label{lem:condition_induced_ii}
    \end{enumerate}
\end{lemma}

\begin{proof}
    \begin{enumerate}[(a)]
    \item By Lemma~\ref{lem:codim1}, it suffices to show that $x$ is a point of codimension one.
    Choose an open affine set $\Spec A = U \subset \X$ such that $x \in U$ and $f \in A$.
    Denote by $\p < A$ the prime ideal corresponding to $x$.
    Since $v$ is non-trivial, the center $x$ is not the generic point of $\X$, and so the codimension is not zero.
    If the dimension of $A/\p$ is at least one, then the inequality
    \[
        \dim{}(A_\p) + \dim{}(A/\p) \leq \dim{}(A) = 2
    \]
    implies
    \[
        \codim(\overline{\{x\}}, \X) = \dim{}(A_\p) \leq 2 - \dim{}(A/\p) \leq 2 - 1 = 1,
    \]
    and hence equals one.
    Thus, it suffices to show that $A/\p$ is not a field.
    If $A/\p$ were a field, it would be a finite extension
    of $k$, but this is impossible since $\bar{f} \in A/\p \subset k(v)$ is transcendental over $k$.
    \item It suffices to show, using (\ref{lem:condition_induced_i}), that $f \in \OO_{\X, x}$.
    Let $P = V(f^{-1})$ and $Z = V(f)$ be the sets of poles and zeros of $f$, respectively.
    Then the rational map $\X \dashrightarrow \PP^1(\OK )$ is obtained by gluing the morphisms
    \[
        \X \setminus P  \to \Spec \OK [t] \subset \PP^1(\OK ), \quad \textrm{and} \quad
        \X \setminus Z \to \Spec \OK [t^{-1}] \subset \PP^1(\OK ),
    \]
    which are induced by
    \begin{center}
        \begin{minipage}{0.25\textwidth}
            \begin{align*}
                \OK [t] &\to \Gamma(\X \setminus P, \OO_{\X}), \\
                t &\mapsto f
            \end{align*}
        \end{minipage}
        \quad and \quad
        \begin{minipage}{0.25\textwidth}
            \begin{align*}
                \OK [t^{-1}] &\to \Gamma(\X \setminus Z, \OO_{\X}), \\
                t^{-1} &\mapsto f^{-1}
            \end{align*}
        \end{minipage}
    \end{center}
    and is defined on $U \coloneqq (\X \setminus P) \cup (\X \setminus Z) = \X \setminus (P \cap Z)$.
    By assumption, $x \in U$. If $x \in \X \setminus P$, the claim follows immediately, as
    \[
        f \in \Gamma(\X \setminus P, \OO_{\X}) \subset \OO_{\X, x}.
    \]
    If $x \in \X \setminus Z$, then
    \[
        f^{-1} \in \Gamma(\X \setminus Z, \OO_{\X}) \subset \OO_{\X, x}.
    \]
    The image of $f^{-1}$ under the local ring homomorphism $\OO_{\X, x} \to \OO_v$ is a unit by assumption,
    hence $f^{-1}$ is also a unit in $\OO_{\X, x}$.
    Consequently $f \in \OO_{\X, x}$.
    \end{enumerate}
\end{proof}

We can now finish the proof of Proposition \ref{prop:constructing_models} (\ref{prop:constructing_models_i}).

\begin{proof}[Proof of Proposition \ref{prop:constructing_models} (\ref{prop:constructing_models_i})]
    Let $f \in \OO_v^\times \subset F_X$ be such that the reduction $\bar{f} \in k(v)$ is transcendental over $k$.
    We also write $f \colon \X \dashrightarrow \PP^1(\OK )$ for the rational map induced by $f$, and let $U$ be its domain of definition.
    Define $\Gamma_f \subset \X \times_{\OK } \PP^1(\OK )$ to be the graph of $f$.
    Recall that it is defined to be the closure of the graph $f|_U$ in $\X \times_{\OK } \PP^1(\OK ) \supset U \times_{\OK } \PP^1(\OK )$.
    This is almost a model of $X$, as it satisfies all the properties of a model except that it might not be normal.
    \begin{center}
        % https://tikzcd.yichuanshen.de/#N4Igdg9gJgpgziAXAbVABwnAlgFyxMJZARgBpiBdUkANwEMAbAVxiRAB12BxOgW17oB9AGYgAvqXSZc+QigAM5KrUYs2nABoACTnl7xBwTgHljggNJid7AAo2AesUMmzl8ZJAZseAkQBMStT0zKyIHLYOTkbsphZi7lLeskRkfsrBamGaCZ7SPnIkpPLpqqHhGgDk4sowUADm8ESgwgBOELxIiiA4EEhk3XRYDGwAFhAQANY5re2d1D1IfhLNbR2IXQuIAMzLIDNrW-O9iAEgDFhgZVAQODi1IEGlbKLUDHQARjAMNnnJYS1YOojHDTVZIAAsR06YgoYiAA
        \begin{tikzcd}
            & \X' \arrow[d]                                &               \\
            \X \times_{\OK } \PP^1(\OK ) \arrow[r,symbol=\supseteq] & \Gamma_f \arrow[r] \arrow[d] & \PP^1(\OK ) \\
            & \X \arrow[ru, "f"', dotted]                  &
        \end{tikzcd}
    \end{center}
    The first projection $\X \times_{\OK } \PP^1(\OK ) \to \X$ induces a birational morphism $\Gamma_f \to \X$.
    The second projection induces a morphism $\Gamma_f \to \PP^1(\OK )$ which extends $f \colon \Gamma_f \dashrightarrow \PP^1(\OK )$.
    Let $\X' \to \Gamma_f$ denote the normalization of $\Gamma_f$.
    This gives a model of $\X'$ for which $f \colon \X' \to \PP^1(\OK )$ is a morphism.
    By Lemma~\ref{lem:condition_induced}~(\ref{lem:condition_induced_ii}), this implies that $v \in V(\X')$.
\end{proof}


\section{Proof of the Second Part of Proposition~\ref{prop:constructing_models}}\label{sec:proof-of-the-second-part}
We fix a model $\X$ of $X$ and a nonempty subset $V \subset V(\X)$.
Denote by $\EE$ the set of irreducible (integral) components of $\X_s$ that induce the valuations in $V(\X)\setminus V$.

\begin{definition}
    A model $\X'$ together with a modification $\varphi \colon \X \to \X'$ such that for every integral vertical curve $E$ on $\X_s$,
    the set $\varphi(E)$ is a point if and only if $E \in \EE$ is called a \textit{contraction} of the $E \in \EE$.
\end{definition}

It can be shown that if a contraction exists, it is unique up to unique isomorphism.
Our goal now is to show that a contraction of $\EE$ exists.
Once this is done, the proof is complete, since for the contraction $\varphi \colon \X \to \X'$, we have by construction $V(\X') = V$.

\begin{lemma}[{{\cite[Lemma 8.3.29, Lemma 8.3.3]{Liu}}}] \label{lem:criterion_point}
    Let $\LL$ be an invertible sheaf on $\X$ generated by global sections $s_0, \dots, s_n$.
    Consider the morphism $f \colon \X \to \PP^n(\OK )$ associated to these sections.
    Let $E$ be a vertical curve of $\X_s$ with the reduced induced closed subscheme structure.
    Then $f(E)$ is reduced to a point if and only if $\LL|_E \cong \OO_E$.
\end{lemma}

\begin{proposition}[{{\cite[Proposition 8.3.30]{Liu}}}] \label{prop:exist_contr}
    The following are equivalent.
    \begin{enumerate}[(a)]
    \item The contraction $\varphi \colon \X \to \X'$ of $\EE$ exists. \label{prop:exist_contr_i}
    \item There exists a Cartier divisor $D$ on $\X$ such that $\deg(D|_{\X_g}) > 0$, that $\LL \cong \OO_\X(D)$ is generated
    by its global sections, and that for any integral vertical curve $E$, we have $\LL|_E \simeq \OO_E$ if and only if $E \in \EE$. \label{prop:exist_contr_ii}
    \end{enumerate}
\end{proposition}

\begin{proof}
    We only sketch (\ref{prop:exist_contr_ii}) $\implies$ (\ref{prop:exist_contr_i}).
    For the other implication and details see~\cite[Proposition 8.3.30]{Liu}.

    As $\X_g$ is an integral projective curve over a field, $\LL|_{\X_g}$
    is an ample divisor (\cite[Proposition 7.5.5]{Liu}). Replacing $D$ by a positive multiple (which
    does not change the set $E$), we can suppose that $\LL|_{\X_g}$ is very ample.

    Let $f \colon \X \to \PP^n(\OK )$ be a morphism associated to sections $s_0, \ldots, s_n \in \Gamma(\X, \LL)$ which generate $\LL$.
    Let $Z$ be the closed subset $f(\X) \subset \PP^n(\OK )$
    endowed with the reduced (hence integral) scheme structure.
    Then $f$ induces a dominant morphism $\X \to Z$ that we still denote by $f$.
    This is an isomorphism on the generic fibre because $\LL|_{\X_g}$ is very ample.
    In particular, $f \colon X \to Z$ is birational.
    Denote by $\pi \colon \X' \to Z$ the normalization morphism.
    As $\X$ is normal, $f$ factors uniquely into $\varphi \colon \X \to \X'$ followed by $\pi \colon \X' \to Z$.
    \begin{center}
        % https://tikzcd.yichuanshen.de/#N4Igdg9gJgpgziAXAbVABwnAlgFyxMJZARgBpiBdUkANwEMAbAVxiRAC0QBfU9TXfIRRkADFVqMWbADrSAGgHJuvEBmx4CREeXH1mrRCFlzu4mFADm8IqABmAJwgBbJACZqOCEm0T9bWyDUDHQARjAMAAr8GkIg9lgWABY4ynaOLojuIJ5IZL5ShrL09miJWKkgDs65Hl6IPnoFRtJo5VwUXEA
        \begin{tikzcd}
            & \X' \arrow[d, "\pi"] \\
            \X \arrow[r, "f"'] \arrow[ru, "\varphi"] & Z
        \end{tikzcd}
    \end{center}
    Let $E$ be an integral vertical curve on $\X$. Then we have
    \[
        f(E) \textrm{ is a point} \iff \LL|_E \simeq \mathcal{O}_E \iff E \in \EE
    \]
    where the first equivalence is a consequence of Lemma~\ref{lem:criterion_point} and the second one is true by assumption.
    We claim that $\pi \colon \X' \to Z$ is a finite morphism.
    This then completes the proof because $f(E)$ is reduced to a point if and only if $\varphi(E)$ is a point.
    The morphism $f$ is projective (\cite[Corollary 3.3.32(e)]{Liu}), so $\AA \cong f_*\OO_\X$ is a coherent
    sheaf on $Z$ (\cite[Corollary 5.3.5]{Liu}). Furthermore, $\AA$ is a sheaf of integrally closed algebras because $\X$ is normal.
    As a result, the normalization morphism $\pi \colon \X' \to Z$
    coincides with the canonical morphism $\SpecRel \AA \to Z$ and is a finite morphism (\cite[Exercise 5.1.17]{Liu}).
\end{proof}

For the construction of a suitable Cartier divisor, we furthermore need the following technical lemma.

\begin{lemma}[{{\cite[Lemma 8.3.32]{Liu}}}] \label{lem:global_sec}
    Let $D$ be an effective horizontal Cartier divisor on $\X$, and $\LL = \OO_\X(D)$.
    Then there exists a $m_0 \geq 1$ such that $\LL^{\otimes m}$ is generated by its global sections for every $m \geq m_0$.
\end{lemma}

Proposition~\ref{prop:exist_contr} and Lemma~\ref{lem:global_sec} show that for a contraction morphism to exist,
it suffices that there exists an effective Cartier divisor that meets the closed fibres in a suitable way.
To construct such a divisor, it is important that the ring $\OK $ is complete, which implies it is Henselian.

\begin{definition} %[{{\cite[8.3.33]{Liu}}}, {{\cite[\href{https://stacks.math.columbia.edu/tag/04GG}{Lemma 04GG}]{Stacks}}}]
    Let $A$ be a Noetherian local ring with maximal ideal $\m$ and residue field $k \coloneqq A/\m$.
    We say that $A$ is Henselian if one of the following equivalent conditions is fulfilled:
    \begin{itemize}
        \item Hensel's Lemma holds: For every monic $f \in A[T]$, and every simple root $a_0 \in k$ of $\bar{f} \in k[T]$,
            there exists a lift $a \in A$ (so $\overline{a} = a_0$) such that $f(a) = 0$.
        \item Every finite $A$-algebra is a direct sum of local $A$-algebras.
    \end{itemize}
\end{definition}

\begin{lemma}[{{\cite[8.3.35]{Liu}}}] \label{lem:exist_div}
    Let $x \in \X_s$ be a closed point on the special fiber, then the following assertions hold.
    \begin{enumerate}[(a)]
        \item There exists an irreducible horizontal Weil divisor $\Delta$ containing $x$. \label{lem:exist_div_i}
        \item There exists an effective horizontal Cartier divisor $D$ such that $\X_s \cap \Supp{D} = \{x\}$. \label{lem:exist_div_ii}
    \end{enumerate}
\end{lemma}

\begin{proof}
    \begin{enumerate}[(a)]
        \item Let $\m_x$ be the maximal ideal of $\OO_{\X,x}$ and denote by $\pi$ a uniformising parameter of $\OO_K$.
            Let $ \p_1, \dots, \p_n $ be the prime ideals of $\OO_{\X,x}$ that are minimal among those containing $\pi \OO_{\X,x}$.
            As $\dim{}\OO_{\X,x} = 2$ (\cite[Proposition 8.3.4]{Liu}), Krull's principal ideal theorem implies that $\m_x \not= \p_i$.
            There therefore exists an
            \[
                f \in \m_x \setminus (\cup_i \p_i).
            \]
            Thus $f$ is a regular element and $V(f)$ does not contain any irreducible component of $\Spec \OO_{\X_s, x}$.

            There exists an affine open subset $W \ni x$ and a regular element $g \in \OO_\X(W)$ such that $g_x = f$,
            and that $V(g) \cap W_s = \{x\}$.
            Let $\Delta$ be the closure in $\X$ of an irreducible component of $V(g)$ passing through $x$.
            Then $\Delta$ is an irreducible horizontal divisor passing through $x$. \label{lem:proof:exist_div_i}
        \item Let $E \coloneqq V(g)$ and let $W$ be as in~\ref{lem:proof:exist_div_i}.
            Then $E$ can be seen as an effective Cartier divisor on $W$.
            Moreover, $E \to \Spec \OO_K$ is quasi-projective and quasi-finite because $\dim{}E_s < \dim{}\X_s = 1$.
            By a corollary of Zariski's Main Theorem (\cite[Corollary 4.4.8]{Liu}), there exists an affine open neighborhood
            $U$ of $x$ such that $E \cap U$ is an open subscheme of a scheme $Z$ that is finite over $\Spec \OO_K$.

            \begin{center}
                % https://tikzcd.yichuanshen.de/#N4Igdg9gJgpgziAXAbVABwnAlgFyxMJZABgBpiBdUkANwEMAbAVxiRAFEACAHW4GM6aTgFUQAX1LpMufIRQBGclVqMWbAFrjJIDNjwEiAJiXV6zVohC8Aymhh8e3APJOA+gGlxymFADm8IlAAMwAnCABbJDIQHAgkRRVzNl4cGAAPHBDw4Ag7ME4YcIAjHygsMF8xEGocOiwGNgALCAgAay1gsMjEBNikY0S1SxT0nGAg8twYKrEKMSA
                % https://tikzcd.yichuanshen.de/#N4Igdg9gJgpgziAXAbVABwnAlgFyxMJZAJgBoAGAXVJADcBDAGwFcYkQAtEAX1PU1z5CKAMwVqdJq3YAdGQGU0MAMYACOQHkNAfQDSPPiAzY8BIuXE0GLNohABRdTOX00qgKo8JMKAHN4RKAAZgBOEAC2SBYgOBBIAIxWUrYgcjgwAB44wEFYYLgw3AbBYZGIZDFxiNHW0nZpmTgh4cAQSmCqMOEARj5Qeb5FNDj0WIzsABYQEADWXtxAA
                \begin{tikzcd}
                    E \cap U \arrow[rr, "\textrm{open emb.}", hook] &  & Z \arrow[r, "\text{finite}"] & \Spec \OO_K
                \end{tikzcd}
            \end{center}

            Let $D$ be the connected component of $Z$ containing $x$.
            Since $\OO_{K}$ is a Henselian local ring, $Z$ is a disjoint union of local schemes.
            Hence $D$ is local, and consequently $D \subset E \cap U \subset \X$. %(Tippfehler in Liu)

            Finally, since $D$ is closed in $Z$, it is finite over $\Spec \OO_K$.
            A finite morphism $D \to \Spec \OO_K$ between affine schemes is projective (\cite[Exercise 3.3.22]{Liu}),
            and therefore proper (\cite[Theorem 3.3.30]{Liu}).
            So the morphisms
            \[
                D \to \X \to \Spec \OO_K \quad
                \textrm{and} \quad
                \X \to \OO_K
            \]
            are proper, which implies that $D \to \X$ is also proper (\cite[Proposition 3.3.16(e)]{Liu}). Consequently, $D$ is closed in $\X$ and therefore a Cartier divisor, as required.
    \end{enumerate}
\end{proof}

Now we can finally prove Proposition~\ref{prop:constructing_models}~(\ref{prop:constructing_models_ii}).

\begin{proof}[Proof of Proposition \ref{prop:constructing_models}~(\ref{prop:constructing_models_ii})]
    Let $Z_1, \dots, Z_n$ be the irreducible components of $\X_s$ that induce the valuations in $V \subset V(\X)$.
    By Lemma~\ref{lem:exist_div}~(\ref{lem:exist_div_ii}), there exists an effective horizontal Cartier divisor $D_i$
    such that $\X_s \cap \Supp D_i$ is a point of $Z_i$ that belongs to none of the components of $\EE$.
    Let $D = \sum_{i=1}^n D_i$.
    By virtue of Lemma~\ref{lem:global_sec}, if necessary replacing $D$ by a multiple, we can even suppose that
    $\LL \coloneqq \OO_\X(D)$ is generated by its global sections.

    The theorem then results from Proposition~\ref{prop:exist_contr}.
    Indeed, \(\LL|_E \simeq \OO_E\) if $E \in \EE$ since $E \cap \Supp D = \varnothing$;
    if $E \notin \EE$, then $E$ is some $Z_i$ and $\deg \LL|_E \geq \deg \OO_\X(D_i)|_{Z_i} > 0$, and hence \(\LL|_E \not\simeq \OO_E\).
\end{proof}

\begin{figure}[h!] \label{fig:illustration-of-proof-of-prop-constructing-models}
    \centering
    \tikzfig{figures/fig1}
    \caption{Illustration of the proof of Proposition \ref{prop:constructing_models}~(\ref{prop:constructing_models_ii})}
\end{figure}

\begin{remark}
    The proof of Proposition~\ref{prop:constructing_models}~(\ref{prop:constructing_models_ii}) shows that actually all models
    $\X$ of $X$ are projective.
    Indeed one just has to choose $V = V(\X)$ in the proof to obtain an effective Cartier divisor $D$ which meets
    all the irreducible components of the fibre $\X_s$.
    Since $\X_s$ is a projective curve (\cite[Exercise 7.5.4]{Liu}), $\OO_\X(D)|_{\X_s}$ is ample (\cite[Proposition 7.5.5]{Liu}).
    It follows from \cite[Corollary 5.3.24]{Liu} that $\OO_\X(D)$ is ample.
    So $\X \to \Spec \OO_K$ is projective (\cite[Proposition 5.3.6]{Liu}).
\end{remark}



\chapter{Base change and permanent models}\label{ch:base-change-and-permanent-models}
\chaptermark{Base change and permanent models}

This chapter studies the behaviour of models and their component valuations under finite extensions of the base field. Its main themes are normalized base change, reduced special fibres and permanent models, the elimination of ramification, and illustrative tame and wild examples.

Throughout this chapter, $X$ is a smooth, projective and absolutely irreducible curve over the local field $K$.

\section{Normalized base change}\label{sec:normalized-base-change}
We want to understand the behaviour of models under base change with respect to an extension of the base field.
Usually, $L|K$ will denote a finite field extension.
Since $v_K$ is Henselian, it has a unique extension to $v_L$ to $L$, and then $(L,v_L)$ also satisfies Assumption \ref{ass:v_K}.
Given a model $\X$ of $X$, we ask ourselves whether the base change $\X\times_{\OO_K}\OO_L$ is again a model of $X_L \coloneqq X\times_K L$.
Well, almost: All properties required of a model from Definition \ref{def:model} hold except possibly normality.
We obtain a model of $X_L$, called the {\em normalized base change} of $\X$ relative to $L|K$, by normalizing $\X\times_{\OO_K}\OO_L$.\footnote{
    Here we use the fact that Assumption \ref{ass:v_K} implies that $\OO_K$ is an \emph{excellent} ring.
    It follows that $\X\times_{\OO_K}\OO_L$ is an excellent scheme, and therefore its normalization is a finite morphism.
    See \cite[Chapter 8, in particular Theorem 8.2.39 (d)]{Liu}.
}

\section{Reduced special fibres and permanence}\label{sec:reduced-special-fibres-and-permanence}
We introduce permanent models and examine their relation to reduced special fibres and weakly unramified component valuations.

We say that a model $\X$ is \textit{permanent} if $\X \times_{\OK}\OL$ is itself a model of $X$:

\begin{definition} \label{def:permanent-model}
    A model $\X$ of $X$ is called \textit{permanent}, if $\X \times_{\OK} \OL$ is normal for all finite field extensions $L|K$.
\end{definition}

We have the following theorem about permanence:

\begin{theorem} \label{thm:permanence}
    \begin{enumerate}[(i)]
        \item For a model $\X$ of $X$,  the following three statements are equivalent:
            \begin{enumerate}[(a)]
                \item $\X$ is permanent,
                \item the special fibre $\X_s$ of $\X$ is reduced,
                \item all $v\in V(\X)$ are \emph{weakly unramified} over $v_K$, i.e.\ $v$ and $v_K$ have the same value group.
            \end{enumerate}
        \item Given a model $\X$, there exists a finite extension $L|K$ such that the normalized base change of $\X$ relative to $L|K$ is permanent.
    \end{enumerate}
\end{theorem}

We will prove part (i) below; see Corollary \ref{cor:permanence-part-i}.
In Section \ref{sec:epp-and-abhyanka}, we will prove part (ii) in the special case where all multiplicities of $\X_s$
are prime to the characteristic of the residue field $k$ (Lemma \ref{lem:abhyanka}).
In later chapters, we see proofs of part (ii) in other important special cases.

\begin{remark}
    Since flatness and properness are preserved under base change, a model $\X$ of $X$ is permanent if and only if
    $\X \times_{\OK} \OL$ is a model of $X_L = X \times_K L$.
\end{remark}

\begin{proposition}[{\cite[Proposition 2.3]{ArzdorfWewers}}] \label{prop:reducedness-implies-permanence}
    If $\X_s$ is reduced, then $\X$ is permanent.
\end{proposition}

\begin{proof}
    Let $L|K$ be a finite extension and $l|k$ the extension of the corresponding residue fields.
    The special fibre of $\X' \coloneqq \X_{\OL} = \X \times_{\OK} \OL$ is $\X'_s := \X_s \times_k l$,
    which is reduced since $k$ is perfect by assumption.
    This means $\X'_s$ satisfies condition (i) from Lemma~\ref{lem:conditions-for-normality}. \\
    For condition (ii) let $\xi \in \X'_s$ be any generic point.
    The stalk $\OO_{\X', \xi}$ then is a Noetherian local domain of dimension $1$ and it holds
    \begin{align}
        \OO_{\X'_s, \xi} = \OO_{\X', \xi}/\pi_L\OO_{\X', \xi},
    \end{align}
    where $\pi_L \in \OL$ is a uniformising parameter.
    Since $\X'_s$ is reduced, $\pi_L\OO_{\X',\xi}$ is a maximal ideal, which means that
    condition (ii) from Lemma~\ref{lem:conditions-for-normality} is satisfied as well.
    Thus $\X'$ is normal and hence $\X$ is permanent.
\end{proof}

\begin{proposition}[{\cite[Corollary 2.2 (ii)]{ArzdorfWewers}}]\label{prop:criterion-reducedness}
    Let $\X$ be a model of $X$, then the following statements are equivalent.
    \begin{enumerate}[(i)]
        \item $\X_s$ is reduced.
        \item $\OO_{\X_s, \xi}$ is a field for every generic point $\xi \in \X_s$.
        \item Every $v_Z \in V(\X)$ is weakly unramified.
    \end{enumerate}
\end{proposition}

\begin{proof}
    Since $\X$ is normal, there are no embedded points in $\X_s$ by Lemma~\ref{lem:conditions-for-normality}.
    Since $\X_s$ has dimension one, it is reduced if and only if $\X_s$ is reduced in codimension zero
    which in turn is the case if $\OO_{\X_s, \xi}$ is a field for every generic point $\xi \in \X_s$. \\
    For the second equivalence recall that according to Proposition~\ref{prop:multiplicities-coincide},
    for any generic point $\xi\in \X_s$ and $Z = \overline{\lbrace\xi\rbrace}$ it holds
    \begin{align}
        [\Gamma_{v_Z}: \Gamma_K] = \length \OO_{\X_s, \xi}.
    \end{align}
    This implies that $\OO_{\X_s, \xi}$ is a field if and only if $v_Z$ is weakly unramified for all $v_Z \in V(\X)$.
\end{proof}

\begin{corollary}[Partial result] \label{cor:permanence-part-i}
    In the notation of Theorem~\ref{thm:permanence}(i), conditions (b) and (c) are equivalent, and each implies condition (a).
    The implication (a) $\Rightarrow$ (b) remains unresolved in the present text.
\end{corollary}

% AUTHOR WARNING:
\begin{authorwarning}
The implication ``permanent $\Rightarrow$ reduced special fibre'' is not proved by the preceding results and does not follow formally from the definition. A correct proof or a precise reference must be supplied.
\end{authorwarning}

\begin{proof}
    The equivalence (b) $\iff$ (c) follows from Proposition \ref{prop:criterion-reducedness}.
    The implication (b) $\Rightarrow$ (a) is Proposition \ref{prop:reducedness-implies-permanence}.
    The remaining implication (a) $\Rightarrow$ (b) is the unresolved point recorded above.
\end{proof}



\section{Extension of component valuations}\label{sec:extension-component-valuations}
To compare the special fibres before and after normalized base change, we describe how their component valuations are related by extension of valued function fields.

Let $L|K$ be a finite field extension and let $\X$ be a model of $X$.
Denote by $\X' \coloneqq (\X\times_{\OK} \OL)^\sim$ the normalized base change of $\X$ to $\OL$.
Let $\xi \in \X_s$ be a generic point and let $\xi' \in \X_s'$ be a generic point lying over $\xi$.
Since $\X$ (resp. $\X'$) is normal we have that $R \coloneqq \OO_{\X, \xi}$ (resp. $S \coloneqq \OO_{\X', \xi'}$)
is a discrete valuation ring dominating $\OK$ (resp. $\OL$) by
Lemma~\ref{lem:conditions-for-normality} (\ref{lem:conditions-for-normality-ii}), that is
\begin{itemize}
    \item $R \supset \OO_K$ (resp. $S \supset \OL$), and
    \item $\m_R \cap \OK = \pi_K\OK$ (resp. $\m_S \cap \OL = \pi_L\OL$).
\end{itemize}
We arrive at the following diagram of fields.

The function field after base change is the compositum
\begin{align}
    F_{X_L}=F_X\mathbin{\cdot}L.
\end{align}

% AUTHOR WARNING:
\begin{authorwarning}
The text assumes that a single rational subfield $K(x)\subset F_X$ can be chosen such that every component valuation $v_i$ restricts to the Gauß valuation. This simultaneous choice requires proof or a precise reference. For an individual valuation such a parameter is easy to choose, but the common choice for all components is stronger. The diagram below records the intended setup only, subject to this unresolved point.
\end{authorwarning}

\begin{figure}[H]\label{fig:valuations-and-field-extensions}
    \centering
    \resizebox{0.5\textwidth}{!}{
        \begin{tikzcd}[ampersand replacement=\&]
            {}
            \& {(F_{X_L},\{w_{i,1},\dots,w_{i,j_i}\})} \\
            {(F_X,\{v_i\})} \arrow[ru, no head]
            \& {} \\
            {}
            \& {(L(x),v_{{\scriptscriptstyle G}, L})} \arrow[uu, no head]  \\
            {(K(x),v_{{\scriptscriptstyle G}, K})} \arrow[ru, no head] \arrow[uu, no head]
            \& {} \\
            {}
            \& {(L,v_L)} \arrow[uu, no head]  \\
            {(K,v_K)} \arrow[ru, no head] \arrow[uu, no head]
            \& {} \\
        \end{tikzcd}
    }
\end{figure}

Here the $v_i$ are the valuations on $F_X$ corresponding to the generic points of the components of $\X_s$.
$w_{i, j}$ represent the valuations corresponding to the irreducible component of $\X_s'$ whose generic point lies over
the generic point of $Z_i \subset \X_s$ which corresponds to $v_i$.


\section{Eliminating ramification}\label{sec:eliminating-ramification}
We now turn to the existence of a finite base extension after which normalized base change has reduced special fibre. The required extension is obtained by eliminating ramification in the component valuations.

\begin{theorem}[{\cite[Proposition 2.3]{ArzdorfWewers}}]
    Let $\X$ be a model of $X$.
    Then there is a finite extension $L|K$ such that $\X' = (\X \times_{\OK}\OL)^\sim$ has a reduced special fibre.
\end{theorem}

% AUTHOR WARNING:
\begin{authorwarning}
A precise version of Epp's theorem appropriate to the finitely many component valuations must be stated here. In particular, the proof must specify the valued fields to which Epp is applied and the required hypotheses on their ranks and residue fields. It must also justify the passage from one component valuation to finitely many of them, explain why a single finite extension $L/K$ works simultaneously, and show that every extension $w_{i,j}$ becomes weakly unramified over $v_L$.
\end{authorwarning}

\begin{proof}
    By~\cite[Theorem 1.0]{Epp} there exists a field extension $L|K$ such that $w_{i,j}$ is weakly unramified over $v_L$.
    According to Proposition~\ref{prop:criterion-reducedness}, the special fibre of $\X'$ is reduced.
\end{proof}


\section{Epp's theorem and Abhyanka's lemma}\label{sec:epp-and-abhyanka}
The preceding existence result invokes Epp's theorem. In the tamely ramified case, the corresponding effect of a base extension on ramification indices is expressed by Abhyanka's lemma.

\begin{lemma}[Abhyanka]\label{lem:abhyanka}
    Let $L|K$ be a finite extension.
    Let $\lbrace v_1, \dots, v_n \rbrace = V(X)$ denote the extended valuation on the function field $F_X|K$.
    We have that $F_{X_L}=F_X\mathbin{\cdot}L$ and with the extended valuation denoted as $w_{1,1}, \dots, w_{n, j_n}$ respectively.
    Assume that $v_i|v_{{\scriptscriptstyle G}, K}$ are tamely ramified for all $i$.
    If $e_i \coloneqq e(v_i|v_{{\scriptscriptstyle G},K})$ divides
    $e(v_{{\scriptscriptstyle G}, L}, v_{{\scriptscriptstyle G}, K}) = e(v_L|v_K)$ for all $i$, then $w_{i,j} \in V(X_L)$
    is weakly unramified over $v_{{\scriptscriptstyle G}, L}$ for all $i,j$.
\end{lemma}

% AUTHOR WARNING:
\begin{authorwarning}
Both the statement and the proof of this lemma require revision. The notation $V(X)$ should presumably be $V(\X)$, the identity $F_{X_L}=F_X\mathbin{\cdot}L$ must be interpreted in a specified common overfield, and the extensions $w_{i,j}$ must be defined. The assertion $w_{i,j}\in V(X_L)$ uses undefined or inconsistent notation, and the ramification indices rely on the unresolved common choice of $K(x)$ recorded above. The present proof is invalid: the reduction from divisibility of ramification indices to equality is unexplained, the displayed relation between uniformisers is incorrect, and $[\Gamma_w:\Gamma_v]=1$ does not follow merely from an equality of compositum fields.
\end{authorwarning}

\begin{proof}
A correct proof must be supplied; see the warning above.
\end{proof}

\begin{comment}
\begin{proof}[Provisional argument retained for the authors]
    We will use the equivalence of Proposition~\ref{prop:criterion-reducedness} without further mention,
    that is the equivalence of the special fibre being reduced and the valuations being weakly unramified.
    Without loss of generality we may assume the following assertions:
    \begin{itemize}
        \item $n = 1$ and $j_1 = 1$ since the argument can be constructed inductively,
        \item $e \coloneqq e_1 = e(v_1|v_{{\scriptscriptstyle G},K}) = e(v_{{\scriptscriptstyle G},L}|v_{{\scriptscriptstyle G},K})$
            since we assume the residue field to be perfect, and
        \item $K$ contains the $e$-th roots of unity, i.e. $\pi_K = \pi_L^e$, since we can always replace $K$ by a suitable
            Kummer extension without violating the other assumptions and obtain a compatible diagram regarding the valuations.
    \end{itemize}
    Since $F_{X_L}=F_X\mathbin{\cdot}L$ we have that $[\Gamma_{w_1}: \Gamma_{v_1}] = 1$.
    By multiplicativity of the ramification indices we conclude $[\Gamma_{w_1}: \Gamma_{v_L}] = 1$.
\end{proof}
\end{comment}


\section{Tame and wild examples}\label{sec:tame-and-wild-examples}
The following examples illustrate how a suitable finite extension and a change of coordinates can produce a model with reduced special fibre, and why the choice of extension matters in the wildly ramified case.

\begin{example}[Obtaining a reduced special fibre after base change]
    Let $K = \QQ_3$ and consider the projective curve
    \begin{align}
        X: \, y^3 = 3x^2 + 1
    \end{align}
    over $K$.
    The equation defines a model $\X$ over $\OK=\ZZ_3$ (cf. Example \ref{ex:model-of-plane-curve}).
    The reduction $\X_s/\FF_3$ is given by $0 = y^3 - 1 = (y - 1)^3$ which means $\X_s$ is not reduced. \\
    The irreducible component is given by $Z: \, y-1$; denote the corresponding valuation by $v$.
    For $z \coloneqq y - 1$ we have
    \begin{align*}
        3x^2 &= y^3 - 1 \\
            &= (z+1)^3 - 1 \\
            &= z^3 + 3z^2 + 3z,
    \end{align*}
    where the right hand side is an Eisenstein polynomial.
    This gives
    \begin{align}
        v(z)=\frac{1}{3}v_{{\scriptscriptstyle G},K}(3x^2)=\frac{1}{3}\notin\ZZ=\Gamma_K.
    \end{align}
    Let us extend the ground field to $L = \QQ_3[\theta]$ where $\theta^3 - 3 = 0$ and apply the coordinate transformation
    $(x,y) \mapsto (x_1, 1+ \theta y_1)$.
    We then obtain a birational curve
    \begin{align}
        X_1: \, y_1^3 + \theta^2 y_1^2 + \theta y_1 = x_1^2.
    \end{align}
    Since $y_1^3 \equiv x_1^2$ (mod $\theta$), $\X_s$ is reduced.
    If $w$ denotes the extension of $v$ associated with this component after base change, then
    $w(\theta y_1)=w(y-1)=\frac{1}{3}=v_L(\theta)\in\Gamma_L$, and hence $w(y_1)=0$.
\end{example}

\begin{example}[Reduction after wildly ramified field extension]
    Consider the curve
    \begin{align}
        X: \, y^2 = 8x^3 - 2
    \end{align}
    defined over the field $K \coloneqq \QQ_2$.
    This example illustrates that the right choice of the field extension is crucial for having a reduced special fibre,
    since a wild ramified field extension is not unique.
    \begin{enumerate}[(i)]
        \item First let us consider the extension $L_1 = K(\theta_1)$ where $\theta_1^2 - 2 = 0$ and the coordinate transformation
            $(x,y) \mapsto (x_1, 2y_1 - \theta_1 - 2)$.
            We obtain a birational curve
            \begin{align}
                X_1: \, y_1^2 - (\theta_1 + 2)y_1 = 2x_1^3 - (2 + \theta_1)
            \end{align}
            The Newton polygon of the resulting quadratic in $y_1$ gives $v(y_1)=\frac14$ for an extension $v$ of
            $v_{{\scriptscriptstyle G},L_1}$ to $F_{X_1}$.
            Since $\Gamma_{L_1}=\frac12\ZZ$ and $[F_{X_1}:L_1(x_1)]=2$, it follows that $[\Gamma_{F_{X_1}}: \Gamma_{L_1}] = 2 \neq 1$.
        \item Let us consider the field extension $L_2 = K(\theta_2)$ where $\theta_2^2 + 2 = 0$
            and the coordinate transformation $(x,y) \mapsto (x_2, 2\theta_2 y_2 + \theta_2)$.
            We obtain a birational curve
            \begin{align}
                X_2: \, y_2^2 + y_2 = -x_2^3
            \end{align}
            which is given by an equation of AS-form, i.e.\ its special fibre is a reduced curve.
            Thus $[\Gamma_{F_{X_2}}: \Gamma_{L_2}] = 1$.
    \end{enumerate}
\end{example}


\part[Analytic and computational tools]{Analytic and computational tools\\[2em]
\begin{minipage}{0.78\textwidth}
\normalfont\normalsize\raggedright
This part develops analytic and computational tools for describing and working with valuations arising from models. The material currently included, Chapter~\ref{ch:mac-lane-valuations}, is essentially the content of the last talk of the miniseminar from which this book originated. It begins the development with Mac Lane valuations, which provide explicit descriptions in terms of polynomials, augmentations, discs, and discoids. The remainder of Part~II will be added later.
\end{minipage}}

\setcounter{chapter}{4}
\chapter{Mac Lane valuations}\label{ch:mac-lane-valuations}

In his seminal work \cite{MacLane}, S.~Mac Lane considered and essentially solved the following problem.
Let $v_K$ be a nontrivial discrete valuation on a field $K$.
What are the extensions of $v_K$ to an (arbitrary) pseudovaluation $v$ on the polynomial ring $K[x]$?
How can we describe such a pseudovaluation explicitly?
These results have been applied to the problem of constructing models of curves in J.~Rüth's thesis~\cite{Rueth}.
This chapter gives a brief overview of pseudovaluations and key polynomials, augmentations and inductive valuations, minimal augmentation chains, and their relation to discs and discoids.
These constructions provide an explicit counterpart to the geometric valuations used in Part~\ref{part:models-and-valuations}.

Fix an algebraic closure $\overline K$ of $K$ and an extension of $v_K$ to $\overline K$, which we again denote by $v_K$.


\section{Pseudovaluations and key polynomials}\label{sec:pseudovaluations-and-key-polynomials}
We first recall pseudovaluations on $K[x]$ and then single out the Mac Lane pseudovaluations and their key polynomials used in the subsequent constructions.

\begin{definition} \label{def:pseudo-valuation}
    We call $v: K[x] \to \hat{\RR}$ a pseudovaluation if $v(0) = \infty$, $v(1) = 0$, and
    \begin{enumerate}[label=\textup{(\alph*)}]
        \item $v(fg) = v(f) + v(g)$,
        \item $v(f + g) \geq \min\lbrace v(f), v(g) \rbrace$ for all $f,g\in K[x]$.
    \end{enumerate}
\end{definition}

\begin{remark}
    Let $v$ be a pseudovaluation on $K[x]$.
    \begin{enumerate}[label=\textup{(\alph*)}]
        \item The definition implies that $\p_v \coloneqq v^{-1}(\infty) < K[x]$ is a prime ideal
            and that $v$ induces a valuation
            \begin{align}
                \bar{v}: K[x]/\p_v \to \hat{\RR}.
            \end{align}
            If $\p_v = (0)$, then $v=\bar{v}$ is a valuation in the usual sense.
        \item Let $\varphi\in K[x]$ be monic and irreducible, and let
            $L=K[x]/(\varphi)$.
            Then we have a one-to-one correspondence between the valuations $v_L: L \to \hat{\RR}$ extending $v_K$ and
            the pseudovaluations $v: K[x] \to \hat{\RR}$ extending $v_K$ with $v(\varphi)=\infty$.
            It is given by the mutually inverse maps
            \begin{align}
                v_L \mapsto \left( K[x] \twoheadrightarrow L \overset{v_L}{\to} \hat{\RR} \right)
            \end{align}
            and
            \begin{align}
                v \mapsto \left(L=K[x]/(\varphi) \overset{\bar{v}}{\longrightarrow} \hat{\RR}\right).
            \end{align}
    \end{enumerate}
\end{remark}

\begin{definition}
    A pseudovaluation $v: K[x] \to \widehat{\QQ}$ is called a \textit{Mac Lane pseudovaluation} if the following conditions hold.
    \begin{enumerate}[label=\textup{(\alph*)}]
        \setcounter{enumi}{2}
        \item $v|_K = v_K$,
        \item $v(x) \geq 0$,
        \item $v$ is discrete, i.e. the subgroup generated by the finite values,
            \begin{align}
                \Gamma_v\coloneqq\bigl\langle v(f):f\in K[x]\setminus\p_v\bigr\rangle,
            \end{align}
            is equal to $\frac{1}{e_v}\ZZ$ for some positive integer $e_v$,
        \item if $\p_v = (0)$, then the extension $k(v)|k$ has transcendence degree one.
    \end{enumerate}
    We denote the set of all Mac Lane pseudovaluations of $K[x]$ by $V(K[x])^*$ and the subset of Mac Lane valuations of $K[x]$ by $V(K[x]) \subset V(K[x])^*$.
\end{definition}

For the valuations in $V(K[x])$, there is a partial order: we say that $v \leq v'$ if and only if $v(f) \leq v'(f)$ for all $f\in K[x]$.

\begin{example}
    The Gauß valuation $v_{\scriptscriptstyle G} \in V(K[x])^*$ is minimal with respect to this partial order.
\end{example}

\begin{definition}
    Let $v\in V(K[x])$ be a Mac Lane valuation and let $f,g,h \in K[x]$.
    \begin{enumerate}[label=\textup{(\alph*)}]
        \item We say that $f$ and $g$ are \textit{$v$-equivalent} (written $f \sim_v g$) if $v(f-g) > v(f) = v(g)$.
        \item The polynomial $f$ is said to \textit{$v$-divide} $g$ (written $f\, |_v\, g$) if $g \sim_v fh$ for some $h \in K[x]$.
            It is called \textit{$v$-irreducible} if $f\, |_v\, gh$ implies $f\, |_v\, g$ or $f\, |_v\, h$,
            and \textit{$v$-minimal} if $f\, |_v\, g$ implies $\deg f \leq \deg g$.
        \item A polynomial $\phi \in K[x]$ is called a \textit{key polynomial} for $v$ if $\phi$ is monic, integral, $v$-irreducible and $v$-minimal.
    \end{enumerate}
\end{definition}

\begin{remark}
    One immediate consequence of these definitions is that any key polynomial $\phi$ for $v$ is irreducible (as an element of $K[x]$).
\end{remark}

\section{Augmentations and inductive valuations}\label{sec:augmentations-and-inductive-valuations}
Starting from a valuation and a key polynomial, augmentation produces a new pseudovaluation by prescribing a value for that polynomial which is at least its previous value. Iterating this construction leads to inductive valuations and their augmentation chains.

Let $\phi$ be a key polynomial for a valuation $v \in V(K[x])$ and let $\lambda \in \hat{\QQ}$ with $\lambda \geq v(\phi)$.
Then we can define a Mac Lane pseudovaluation
\begin{align} \label{eq:def-augmentation}
    v' = [v, v'(\phi) = \lambda] \in V(K[x])^*
\end{align}
as follows.
Any polynomial $f \in K[x]$ can be written uniquely as
\begin{align}
    f = \sum_i f_i \phi^i
\end{align}
with $\deg f_i < \deg \phi$. Then let
\begin{align}
    v'(f) \coloneqq \min_i(v(f_i) + i\cdot \lambda).
\end{align}

By construction, $v\leq v'$, and $v < v'$ if and only if $v(\phi) < \lambda$.

\begin{definition}
    The Mac Lane pseudovaluation $v'$ defined in~\eqref{eq:def-augmentation} is called the \textit{augmentation} of $v$ with respect to $(\phi,\lambda)$.
    If $v < v'$, $v'$ is called a \textit{proper augmentation}.
\end{definition}

\begin{definition}
    A Mac Lane pseudovaluation $v$ is called \textit{inductive} if there exists a chain of Mac Lane pseudovaluations
    \begin{align}
        v_{\scriptscriptstyle G} = v_0 \leq v_1 \leq \cdots \leq v_n = v,
    \end{align}
    such that $v_i$ is a proper augmentation of $v_{i-1}$ for $i=1,\dots, n$.
    We call $v_0,\dots,v_n$ the \textit{augmentation chain} representing $v$ and write \begin{align}
        v = [v_0,v_1(\phi_1) = \lambda_1,\dots, v_n(\phi_n) = \lambda_n].
    \end{align}
    Here $v_0 = v_{\scriptscriptstyle G}$, $v_n = v$, the polynomial $\phi_i$ is a key polynomial for $v_{i-1}$, $\lambda_i > v_{i-1}(\phi_i)$,
    and $v_i = [v_{i-1}, v_i(\phi_i) = \lambda_i]$.
    Also, $\lambda_i \neq \infty$ for $i < n$.
    We say that the augmentation chain $v_0,\dots,v_n$ is \textit{minimal} if
    \begin{align}
        \deg \phi_1 < \cdots < \deg \phi_n.
    \end{align}
\end{definition}


\section{Minimal augmentation chains}\label{sec:minimal-augmentation-chains}
The central structure result for these constructions describes Mac Lane pseudovaluations by minimal augmentation chains.

The following result combines the existence statements in~\cite[Theorem 4.31 and Corollary 4.67]{Rueth} with the uniqueness criterion in~\cite[Theorem 4.33]{Rueth}.

\begin{theorem}[Minimal augmentation chains]\label{thm:minimal-augmentation-chains}
    \begin{enumerate}[label=\textup{(\roman*)}]
        \item Every $v\in V(K[x])$ is inductive and can be represented by a minimal augmentation chain.
        \item If $K$ is complete, then every $v\in V(K[x])^*$ is inductive and can be represented by a minimal augmentation chain.
        \item Let
        \begin{align*}
            v_n&=[v_0,v_1(\phi_1)=\lambda_1,\dots,v_n(\phi_n)=\lambda_n],\\
            w_m&=[v_0,w_1(\psi_1)=\mu_1,\dots,w_m(\psi_m)=\mu_m]
        \end{align*}
        be inductive pseudovaluations represented by minimal augmentation chains.
        Then $v_n=w_m$ if and only if $m=n$ and, for every $i=1,\dots,n$,
        \begin{enumerate}[label=\textup{(\alph*)}]
            \item $\lambda_i=\mu_i$,
            \item $\deg\phi_i=\deg\psi_i$, and
            \item $v_{i-1}(\phi_i-\psi_i)\geq\lambda_i$.
        \end{enumerate}
    \end{enumerate}
\end{theorem}

\section{Discs, discoids, and valuations}\label{sec:discs-discoids-and-valuations}
We conclude by introducing discs and discoids in $\overline K$, which give a geometric description of subsets defined by valuation conditions on polynomials.

\begin{definition}
    \begin{enumerate}[label=\textup{(\alph*)}]
        \item Let $\alpha \in \overline K$ and let $\lambda \in \hat{\QQ}$.
            Then
                \begin{align}
                    D(\alpha, \lambda) \coloneqq \lbrace x \in \overline K:\, v_K(x - \alpha) \geq \lambda\rbrace
                \end{align}
            is the \textit{(closed) disc} with centre $\alpha$ and radius $\lambda$.
        \item Let $\phi \in K[x]$ be a monic irreducible polynomial and let $\lambda \in \hat{\QQ}$.
            The set
            \begin{align}
                D(\phi, \lambda) \coloneqq \lbrace x \in \overline K:\, v_K(\phi(x)) \geq \lambda \rbrace
            \end{align}
            is called the \textit{discoid} with centre $\phi$ and radius $\lambda$.
    \end{enumerate}
\end{definition}

\begin{remark}
    For the following description, assume that $K$ is complete, as in the cited source.
    Any disc with centre in $K$ is a discoid.
    Conversely, if $\phi = x - \alpha$ then $D(\phi, \lambda)$ coincides with the disc $D(\alpha, \lambda)$.
    More generally, let $\phi\in K[x]$ be monic and irreducible.
    Let $\alpha_1,\dots,\alpha_n\in\overline K$ be the roots of $\phi \in K[x]$, listed with multiplicity, so that $n=\deg\phi$,
    and let $\mu_1,\dots,\mu_n \in \hat{\QQ}$ be minimal\footnote{
        Our definition of discs implies that discs with large radius are smaller than those with small radius.
        For instance, the disc with centre $0\in K$ and radius $0$ is relatively large, as it consists of all integral elements.
    } such that $D(\alpha_i,\mu_i) \subset D(\phi,\lambda)$.
    Then
    \begin{align}
        D(\phi,\lambda) = \bigcup_{i=1}^n D(\alpha_i,\mu_i),
    \end{align}
    where the discs in this union need not be pairwise distinct; see \cite[Lemma 4.43]{Rueth}.
\end{remark}


\backmatter

\bibliographystyle{alpha}
\bibliography{references}

\end{document}
